\chapter{Implémentation}

\section{Phase exploratoire}
\index{Phase exploratoire}

\subsection{Introduction}
\index{introduction}
La première partie du projet a consisté à élaborer une démonstration, afin de pouvoir correctement définir le cahier des charges selon des besoins réels. Cette phase a été essentielle pour comprendre les enjeux du projet et pour orienter les choix technologiques et méthodologiques qui ont suivi.

Il a été décidé que le mieux pour cette phase était de créer un programme de démonstration, qui simule une machine industrielle. Ce programme représente un HMI (Human-Machine Interface) d'une machine industrielle. Il ne s'agit que d'une facade, car il n'y a aucune machine physique ou virtuelle derrière. Il s'agit seulement une interface graphique. Son élaboration a été faite un peu à l'aveugle, en ajoutant des fonctionnalités au fur et à mesure de l'avancement de cette phase exploratoire.

L'objectif était de voir quelles sont les fonctionnalités qui sont nécessaires pour la version finale, lequelles sont utiles et lesquelles sont superflues.

\subsection{Front-end}
\index{front-end}
% Il faudrait expliquer le fonctionnement du front-end, avec des captures d'écran, etc...

\subsection{Back-end}
\index{back-end}
% Il faudrait expliquer le fonctionnement du back-end, avec des éxtraits de code, etc...

\subsection{MCP}
\index{MCP}
% Il faudrait expliquer le fonctionnement du MCP, avec des éxtraits de code, etc...

\subsection{RAG}
\index{RAG}
% Il faudrait expliquer le fonctionnement du RAG, avec des éxtraits de code, etc...

\subsection{Hebergement de LLM}
\index{hebergement}
% Il faudrait expliquer comment le LLM est hébergé, avec des éxtraits de code, etc...

\subsection{Tests des biais des LLM}
\index{biais}
% Il faudrait expliquer comment les biais des LLM ont été testés, etc...

\subsection{Conclusion}
\index{conclusion}
% Il faudrait conclure sur cette démonstration et qu'est-ce qu'elle a apporté
